\chapter{Notes for later: Not a real chapter. Ignore}
Add content in chapters/02.tex

Bebop players in the 40s started playing sharp 4 or 11 over a major 4th chord: typically a 4th.

Humpty Dumpty made words mean just what he wanted them to mean. He expanded the semantic Overton window.

The grammaticalization of NPIs (Israel2001), or grammaticalization in general

\begin{quote}
    Perhaps the most surprising asymmetry in the expression of polarity is the phenomenon of polarity sensitivity—the tendency for certain forms, polarity items, to be distributed unevenly across negative and affirmative contexts. Although the details vary from language to language, polarity items seem likely to occur in every human language. In many there are literally hundreds of such forms. This is actually rather odd: if one were to devise an artificial language, the idea of including forms which are systematically excluded from certain sentences might seem a perverse extravagance; however, it is an extravagance which natural languages commonly indulge. (\cite[11]{Israel2004})
\end{quote}
\begin{quote}
    polarity items tend to resist purely structural explanations, and pragmatics often plays a role in explaining the details of their sensitivities. Ultimately, I suggest, their grammaticality is a matter not of grammar alone, but depends crucially on their rhetorical fit with the contexts in which they occur.\cite[13]{Israel2004}
\end{quote}

\begin{quote}
    abe: uhhuh I heard of that in my life. \cite[5]{Israel2011}
\end{quote}

\begin{quote}
    The “grammaticality” of a polarity item in a linguistic context is thus a function not just of the sentence in which it occurs, but also of the utterance in which it is used.\cite[9]{Israel2011}
\end{quote}

\begin{quote}
    (ll) (i) I called the man who wrote the book that you told me about up \cite[11]{Chomsky1965}
\end{quote}

\begin{quote}
    Ladusaw (1979) first proposed that licensing is a matter of logical semantics, it has been widely noted that the relevant infer- ences are subject to pragmatic constraints: they can be triggered by certain con- versational implicatures (Linebarger 1980, 1987, 1991; Israel 1996), and they may be blocked if certain presuppositions are not held constant (Heim 1984; Kadmon \& Landman 1993; von Fintel 1999; Horn 2002). \cite[61]{Israel2011}
\end{quote}

\begin{quote}
    I contend that polarity contexts are defined not just by the logical, semantic, or syntactic properties of constructions, but also, and cru- cially, by the pragmatic properties which determine how an expressed propos- ition is construed in context. The relevant level of analysis here is thus neither purely linguistic nor purely conceptual, but a little of both: it is the level at which a speech act is conceptualized in an act of communication, the mental space in which a proposition is expressed and interpreted. \cite[62]{Israel2011}
\end{quote}

The development of NPIs.

\section{Over-imitation}
The thing

Grammaticality is not a unified concept; and although some of the ideas associated with grammaticality may show one or more of these characteristics, none of them is essential to it, but at most only contingent. (paraphrased from John Joseph)

Signed languages can assign referents to physical spaces, allowing 


It's only a bit of an exaggeration to say that the idea of humans having a universal grammar is a bit like the idea that the earth has a universal blueprint for life. In both cases, once things get started, there are evolutionary pressures that favour particular structures. The fact that the earth has gravity, water, sunlight, and carbon dictates what life may evolve. The fact that humans have ears, eyes, mouths, hands, and a brain with certain attentional proclivities similarly dictates what languages may evolve. In both cases, many unexplored possibilities exist, and this is in part because there is no simple path between the realized and unrealized forms.

The common ancestor of vertebrates and animals with exoskeletons dates back to the early Cambrian period or possibly the late Pre-Cambrian, over 500 million years ago. This ancestor would have been a simple, soft-bodied organism. From this point, the evolutionary paths diverged, with one leading to the development of invertebrates with exoskeletons (e.g., insects) and the other leading to chordates, which eventually gave rise to vertebrates. Having made this choice, it's very unlikely that any insects will ever trade their exoskeleton for bones or that any mammals will make the opposite exchange.

Similarly, English and the Germanic languages have evolved prepositions while Persian and the Indo-Iranian languages have maintained the postpositions of Proto-Indo-European. At this point, it's very unlikely that either language will flip ``position'' direction.

But we don't actually have a grammar in our DNA, in our brains. We have constraints, preferences, and tendencies that shape the way language can evolve, the way it's likely to evolve. When we say something is ungrammatical, we're not saying that it violates some fundamental rule in our human nature. Rather, we're saying something much simpler, something much more like, "I can't see why you would say it like that."

``About 60 years after Saussure, the philosopher David Lewis characterised convention as a relationship between a type of agent and a behaviour that is stable, collectively known, and of course, arbitrary."

``So just by dint of acting as if it already was common knowledge, it now is and will remain in the future. That is the funny thing: it can be willed into existence by believing it was already there.
\dots
Each time someone speaks to us, the choices of words, their intonation, the idioms they use, and so on are presupposing a language, which we accommodate. But everyone else is doing the same, accommodating the language we produce.

In short, the continual attempt, by each of us, to determine our language’s nature is the force that shapes it over time. Everyone acts as if the rules are already common knowledge, and then – the strangest part! – they gradually try to work out what this common knowledge is.''https://aeon.co/essays/are-the-exact-words-of-a-language-arbitrary-or-necessary

\section{Intentionally}

See discussion in \citet[315--]{Hurford2007}.

https://www.agentquery.com/search.aspx

``In Chapter 6, I discuss one odd exception to the claim that most natural language has local dependencies: legal language, sometimes called Legalese. It turns out that Legalese has unusually long dependencies'' (Gibson, 2024 \textit{Syntax})